\hypertarget{chapter-19-debugging-like-a-superuser}{%
\chapter{Debugging Like a Superuser}\label{chapter-19-debugging-like-a-superuser}}

The defining characteristic of a superuser is not that they never encounter errors, but that they know how to interpret them. Most LaTeX users respond to a failed compilation by changing random lines until the error disappears, or worse, by copying fragments from internet forums without understanding why they work. This treats debugging as guesswork. A superuser instead regards the compiler as an instrument reporting the current state of the document, and treats the task not as eliminating an error by any means available, but as understanding why the document entered an inadmissible state and what minimal intervention restores it.

\hypertarget{reading-the-log-file}{%
\section{Reading the Log File}\label{reading-the-log-file}}

The first habit worth cultivating is reading the \texttt{.log} file directly rather than relying on an editor's abbreviated error summary. Every compilation produces a detailed account of package loading, macro expansion, font selection, and page breaking, and the message an editor surfaces is often only the final, downstream consequence of a problem that originated much earlier. A missing brace, for instance, frequently does not produce an error at the point it is missing; TeX continues parsing, absorbing subsequent text into whatever group it believes is still open, until it can no longer maintain a consistent interpretation of the input, at which point an error surfaces hundreds of lines later, in a location that has nothing to do with the actual mistake. The \texttt{.log} file's complete execution history, read from the point just before the reported error backward, is what actually identifies the original cause; the editor's summary, taken alone, frequently points a debugging effort in the wrong direction entirely.

\hypertarget{expansion-errors-versus-runtime-errors}{%
\section{Expansion Errors versus Runtime Errors}\label{expansion-errors-versus-runtime-errors}}

Understanding TeX's expansion model, covered in Part I and Part III of this book, matters directly here, because many apparent runtime failures are in fact expansion failures with a runtime-looking symptom. A document does not fail because some command executed incorrectly in the way a runtime error would in a conventional language; it fails because the wrong sequence of tokens was produced during expansion, before anything resembling execution began, and the resulting malformed token sequence only produces a visible error once something further downstream tries and fails to make sense of it.

\texttt{\textbackslash{}tracingmacros} and \texttt{\textbackslash{}tracingcommands}, set to a nonzero value and written to the log with the aid of \texttt{\textbackslash{}tracingonline} during a targeted debugging session, reveal this hidden expansion process directly: every macro expansion is recorded as it happens, showing exactly what a macro's arguments were and what it expanded to, which makes it possible to observe precisely where an unexpected token sequence first appeared rather than inferring it indirectly from a downstream symptom. \texttt{\textbackslash{}show} and \texttt{\textbackslash{}showthe}, used more surgically at a single point in a document, report a macro's current definition or a register's current value without needing to enable full tracing, and are usually the faster tool for a specific, localized question (``what does this macro currently expand to'') rather than a broad investigation. These tools are verbose, and reading their output takes practice, but they are the direct evidence Part I's model predicts should exist, and using them is the difference between confirming a hypothesis about where a problem originates and merely guessing at it.

\hypertarget{profiling-compilation-performance}{%
\section{Profiling Compilation Performance}\label{profiling-compilation-performance}}

Performance problems deserve the same systematic treatment as correctness problems. A large monograph, with thousands of pages, hundreds of figures, and an extensive bibliography, frequently compiles slowly enough that the temptation is to accept the slowness as inevitable rather than to investigate it. It is rarely uniform: external graphics, TikZ figures recomputed from scratch on every build, syntax-highlighted listings processed through \texttt{minted}'s external Pygments call, glossary and index generation, and bibliography processing each contribute a measurable, separately identifiable cost, and \texttt{texloganalyser} or a comparable log-timing tool, alongside simple wall-clock timing of individual build stages through the Makefile targets described in Chapter 5, can identify which of these actually dominates a given project's build time. A project whose TikZ figures are the primary cost benefits from caching compiled figures between builds and only regenerating the ones that changed; a project whose bibliography processing dominates benefits from investigating whether Biber's own caching is being invalidated unnecessarily. Guessing at which stage is slow and optimizing it without measurement frequently targets the wrong stage entirely, spending effort on a part of the build that was never the actual bottleneck.

\hypertarget{a-philosophy-of-repair-the-nearest-admissible-state}{%
\section{A Philosophy of Repair: The Nearest Admissible State}\label{a-philosophy-of-repair-the-nearest-admissible-state}}

The overwhelming majority of compilation failures arise from a small, local inconsistency: an unmatched delimiter, a malformed macro argument, an incorrectly nested environment, a missing package. The document as a whole remains largely valid; only some small part of it has drifted into a state the compiler cannot process. The objective, correctly understood, is to restore the nearest state the compiler can accept, not to reconstruct the project from memory or rewrite substantially more than what actually broke.

This can be stated a little more precisely without turning it into a heavier formal exercise than it needs to be. Treat the space of possible document states as everything the source files could contain, and treat a state as admissible if it compiles to a valid, intended output; a failed build is a document that has moved, through some edit, from an admissible state into an inadmissible one. Debugging, under this framing, is the search for a repair, a small edit, ideally the smallest one that actually resolves the failure, that returns the document to an admissible state again. This is worth stating explicitly because it reframes what ``success'' in debugging looks like: the goal is not any edit that makes the error message disappear, since a sufficiently large or careless edit can often suppress an error while introducing a different, sometimes less visible, problem elsewhere; the goal is the minimal repair, the smallest change consistent with removing the actual inconsistency, because that is the change most likely to restore the document to what it was actually meant to say, rather than to some other admissible state that happens to compile but no longer matches the author's intent.

\hypertarget{version-control-as-a-search-procedure}{%
\section{Version Control as a Search Procedure}\label{version-control-as-a-search-procedure}}

This framing turns directly into a practical technique once a project follows Chapter 4's version control discipline: if the document compiled successfully at some earlier commit and fails now, the search for a repair is constrained to whatever changed between that commit and the present, rather than spanning the entire manuscript. \texttt{git\ bisect}, applied to a build failure exactly as it would be applied to a software regression, automates this search: given a known-good commit and a known-bad one, it checks out successive commits between them, compiling at each, until it identifies the exact commit that introduced the failure, at which point the diff for that single commit is almost always small enough to make the actual cause immediately visible.

Small, logically coherent commits, as Chapter 4 already recommended for reasons of readable history, pay an additional dividend here: a bisection that lands on a commit touching a single chapter file is far more informative than one landing on a commit that happened to touch a dozen unrelated files at once, since the former narrows the search to a specific paragraph while the latter still leaves a meaningful amount of investigation to do even after the offending commit is identified. Compilation, treated this way, stops being a final hurdle checked only when a chapter feels finished and becomes a continuous verification process, ideally run on every commit through the CI setup from Chapter 15, so that a failure is caught and localized to a single small commit immediately rather than discovered much later, once dozens of subsequent changes have accumulated on top of it and the search space has grown accordingly.

\hypertarget{isolating-package-and-workflow-changes}{%
\section{Isolating Package and Workflow Changes}\label{isolating-package-and-workflow-changes}}

The same discipline extends to package upgrades and workflow changes, not only to content edits. Introducing a new document class, switching bibliography systems, or upgrading a graphics package should be done in its own isolated commit, ideally its own isolated test build, so that any resulting incompatibility has an unambiguous origin. Changing several such things simultaneously, upgrading a package and switching engines and revising the class file all in one pass, makes diagnosis considerably harder than it needs to be, since a resulting failure could plausibly trace to any of the changes, and confirming which one is actually responsible requires undoing them one at a time regardless of whether they were introduced one at a time or all at once. Introducing them one at a time in the first place avoids that reconstruction entirely.

Debugging in LaTeX, treated with this discipline throughout, is ultimately an exercise in preserving continuity. A project has a most recent working state, and the purpose of debugging is to find the smallest repair that reconnects the current document to that state while keeping whatever genuine improvement motivated the change that broke it. This mirrors good software engineering generally: a reliable project is maintained not through heroic rewrites undertaken in a moment of frustration with a broken build, but through careful observation of what the compiler is actually reporting, incremental modification narrow enough to be individually verifiable, and disciplined repair aimed at the smallest change that restores a working document.

Part VII closes here. Part VIII turns from individual documents to the corpus-level discipline of writing and maintaining book-length works over years, and to the broader ecosystem of tools a superuser's daily workflow depends on.
